\documentclass[11pt]{article}

% --- UNIVERSAL PREAMBLE BLOCK ---
\usepackage[a4paper, left=1.06cm, top=1.7cm, right=1.06cm, bottom=1.5cm]{geometry}
\usepackage{fontspec}

\usepackage[bidi=basic, provide=*]{babel}

\babelprovide[main, import, onchar=ids fonts]{spanish}
\babelprovide[import, onchar=ids fonts]{english}

% Set default/Latin font to Sans Serif in the main (rm) slot
\babelfont{rm}{Noto Sans}

% Add because we replaced the legacy list behavior
\usepackage{enumitem}
\setlist[itemize]{label=-}
% --- END UNIVERSAL PREAMBLE BLOCK ---

\usepackage{hyperref}
\usepackage{etoolbox}
\setlength{\parindent}{0pt}
\providecommand{\CVProfile}{ai-engineer}
\providecommand{\IncludeSummary}{1}
\newcommand{\ProfileSwitch}[3]{%
    % #1 academic, #2 fullstack-engineer, #3 ai-engineer
    \ifstrequal{\CVProfile}{academic}{#1}{%
        \ifstrequal{\CVProfile}{fullstack-engineer}{#2}{#3}%
    }%
}
\newcommand{\ProfileSummary}{\ProfileSwitch{
% academic
M.C. en Ciencias de la Computacion con experiencia docente y de mentoria, enfocado en educacion de IA/ML y proyectos aplicados que conectan investigacion y producto.
}{
% fullstack-engineer
Ingeniero full-stack con experiencia en construccion y entrega de plataformas en produccion, integrando backend, frontend e infraestructura cloud.
}{
% ai-engineer
Ingeniero de IA con experiencia full-stack en construccion y entrega de sistemas en produccion con Python, Ruby on Rails, React y plataformas cloud.
}}

\begin{document}
\begin{center}
    \textbf{\LARGE Miguel Angel Bravo Vidales}\\ 
    \vspace{4pt}
    \hrulefill
\end{center}

\begin{center}
    San Luis Potosí, SLP, México \textbullet \ x@mackaber.me \textbullet \ +52 444 444 7397 \textbullet \ linkedin.com/in/mackaber \textbullet \ github.com/mackaber
\end{center}

\ifnum\IncludeSummary=1
\vspace{6pt}
\begin{center}
    \textbf{Resumen}
\end{center}
\ProfileSummary
\fi

\vspace{0.5pt}

\begin{center}
    \textbf{Educación}
\end{center}
\textbf{Tecnológico de Monterrey (ITESM)} \hfill Monterrey, NL

\ProfileSwitch{
% academic
Maestría en Ciencias Computacionales \hfill 2018 -- 2019
}{
% fullstack-engineer
Maestría en Ciencias de la Computación \hfill Dic 2020
}{
% ai-engineer
Maestría en Ciencias de la Computación (Enfoque en Inteligencia Artificial y Aprendizaje Automático) \hfill Dic 2020
}

\vspace{10pt}

\textbf{Tecnológico de Monterrey (ITESM)} \hfill San Luis Potosí, SLP

\ProfileSwitch{
% academic
Ingeniería en Tecnologías de Información y Comunicaciones (ITIC) \hfill 2009 -- 2013
\begin{itemize}[noitemsep, topsep=2pt, partopsep=0pt, parsep=0pt]
    \item Premio DAE ``Al desarrollo estudiantil''
    \item Premio ``CENEVAL al Desempeño de Excelencia-EGEL'' (2014)
\end{itemize}
}{
% fullstack-engineer
Ingeniería en Tecnologías de Información y Comunicaciones \hfill Dic 2013
}{
% ai-engineer
Ingeniería en Tecnologías de Información y Comunicaciones \hfill Dic 2013
}

\vspace{12pt}

\begin{center}
    \textbf{Experiencia}
\end{center}

\textbf{Tecnológico de Monterrey (ITESM)} \hfill San Luis Potosí, SLP

\textbf{Profesor de Preparatoria} \hfill Agosto 2025 -- Presente
\begin{itemize}[noitemsep, topsep=0pt, partopsep=0pt, parsep=0pt]
    \item Impartí clases de desarrollo de software a estudiantes de preparatoria, cubriendo fundamentos de programación, construcción de lógica y diseño práctico de aplicaciones.
\end{itemize}

\vspace{8pt}

\textbf{Encora, Inc.} \hfill Remoto

\textbf{Instructor de AWS re/Start} \hfill Febrero 2025 -- Junio 2025
\begin{itemize}[noitemsep, topsep=0pt, partopsep=0pt, parsep=0pt]
    \item Impartí capacitación integral en cómputo en la nube a un grupo de estudiantes, cubriendo servicios principales de AWS, redes, seguridad y fundamentos de Linux/Python.
    \item Guié a los estudiantes en laboratorios prácticos para prepararlos para la certificación AWS Certified Cloud Practitioner y su transición a roles de entrada en cómputo en la nube.
\end{itemize}

\vspace{8pt}

\textbf{Spark Tech Staff} \hfill Enero 2025 -- Presente
\begin{itemize}[noitemsep, topsep=0pt, partopsep=0pt, parsep=0pt]
    \item Mentoricé a nuevas contrataciones en estructuras de datos y algoritmos mediante el programa interno Algorithm School Program, acelerando su preparación para asignación con clientes.
\end{itemize}

\vspace{8pt}

\textbf{Desarrollador Full Stack} \hfill Enero 2021 -- 2024
\begin{itemize}[noitemsep, topsep=0pt, partopsep=0pt, parsep=0pt]
    \item Construí herramientas de desarrollo impulsadas por IA usando Python, LangChain, Streamlit y la API de OpenAI, automatizando la generación de aplicaciones plantilla en React y React Native (Expo).
    \item Mantuve y optimicé una aplicación heredada de lealtad para viajes (Arrivia), implementando monitoreo de recursos con Azure AppInsights para identificar y resolver tareas de alto consumo.
    \item Desarrollé pipelines robustos de procesamiento de datos para una plataforma de telesalud (MDLive), procesando expedientes de pacientes y construyendo sistemas de incentivos para médicos con Ruby on Rails, Redis y RabbitMQ.
\end{itemize}

\vspace{8pt}

\textbf{DEV.f} \hfill Remoto

\textbf{Profesor de Desarrollo Web} \hfill Agosto 2021 -- Presente
\begin{itemize}[noitemsep, topsep=0pt, partopsep=0pt, parsep=0pt]
    \item Impartí bootcamps de desarrollo web full-stack de 6 meses cubriendo HTML, CSS, JavaScript, Node.js y React.
    \item Guié a estudiantes sin experiencia previa en desarrollo de software mediante proyectos prácticos para construir portafolios listos para la industria.
\end{itemize}

\vspace{8pt}

\textbf{Continental} \hfill San Luis Potosí, SLP

\textbf{Desarrollador Full Stack (Freelancer)} \hfill Enero 2014 -- Diciembre 2017
\begin{itemize}[noitemsep, topsep=0pt, partopsep=0pt, parsep=0pt]
    \item Desarrollé y di mantenimiento a la aplicación Android para una plataforma de administración de transportadores.
    \item Construí un servicio de integración de escritorio en Ruby para conectar nuevas interfaces con aplicaciones legadas existentes.
\end{itemize}

\vspace{8pt}

\textbf{Winero} \hfill San Luis Potosí, SLP

\textbf{Cofundador, CTO \& Desarrollador Full Stack} \hfill Julio 2013 -- Diciembre 2016
\begin{itemize}[noitemsep, topsep=0pt, partopsep=0pt, parsep=0pt]
    \item Cofundé una plataforma de programas de lealtad, diseñando la arquitectura de datos, el sistema backend (Ruby on Rails, PostgreSQL) y el panel de analítica de rendimiento.
\end{itemize}

\vspace{12pt}

\ProfileSwitch{
% academic
\begin{center}
    \textbf{Conferencias y Cursos}
\end{center}
\begin{itemize}[noitemsep, topsep=0pt, partopsep=0pt, parsep=0pt]
    \item Taller de Hacking, Tecnológico de Monterrey Campus San Luis Potosí (2012)
    \item Conferencia de Hacking (CongresoExit), Tecnológico de Monterrey Campus San Luis Potosí (2014)
    \item Conferencia de Hacking de Aplicaciones Web, Tecnológico de Monterrey Campus Monterrey, SISCTI 40 (2015)
    \item Conferencia de Hacking de Aplicaciones Web, Instituto Tecnológico de Piedras Negras (2015)
    \item Tuna Hack, hackathon de programación (2015, 2017)
    \item Technovation, hackathon de programación enfocado en niñas (2016, 2017)
    \item TunaTech, cumbre de tecnología y negocios (2017, 2019)
    \item Talent Hackathon (2019 -- 2023)
    \item Wolfram Summer School (2019)
\end{itemize}

\vspace{12pt}

\begin{center}
    \textbf{Participación Académica y Comunitaria}
\end{center}
\begin{itemize}[noitemsep, topsep=0pt, partopsep=0pt, parsep=0pt]
    \item Cofundé ``Comunidad Código,'' una asociación de programadores en San Luis Potosí.
    \item Colaboré con el equipo organizador de la asociación ``DesarrolladorSLP''.
    \item Participé en más de 17 hackathons: 10 como asistente, 7 como mentor, 1 como organizador y ganador en 2.
    \item Apoyé la creación de Tuna Valley, red de comunidades locales en San Luis Potosí.
\end{itemize}
}{
% fullstack-engineer
\begin{center}
    \textbf{Liderazgo y Actividades}
\end{center}

\textbf{Comunidades y Eventos Tech} \hfill México

\textbf{Speaker, Mentor y Cofundador} \hfill 2013 -- Presente
\begin{itemize}[noitemsep, topsep=0pt, partopsep=0pt, parsep=0pt]
    \item Participé en más de 20 hackathons (ganando 2 y mentoreando 10), incluyendo Startup Weekend, AngelHack y Talent Hackathon.
    \item Impartí más de 15 charlas y talleres sobre ingeniería de software y desarrollo de plataformas en eventos tecnológicos como TalentLand, Nerdearla y NerdyTalks.
    \item Cofundé ``Comunidad Código,'' una asociación local de desarrolladores, y contribuí activamente a la iniciativa Tuna Valley.
\end{itemize}
}{
% ai-engineer
\begin{center}
    \textbf{Liderazgo y Actividades}
\end{center}

\textbf{Comunidades y Eventos Tech} \hfill México

\textbf{Speaker, Mentor y Cofundador} \hfill 2013 -- Presente
\begin{itemize}[noitemsep, topsep=0pt, partopsep=0pt, parsep=0pt]
    \item Participé en más de 20 hackathons (ganando 2 y mentoreando 10), incluyendo Startup Weekend, AngelHack y Talent Hackathon.
    \item Impartí más de 15 charlas y talleres sobre ingeniería de software e IA en eventos tecnológicos como TalentLand, Nerdearla y NerdyTalks.
    \item Cofundé ``Comunidad Código,'' una asociación local de desarrolladores, y contribuí activamente a la iniciativa Tuna Valley.
\end{itemize}
}

\vspace{12pt}

\begin{center}
    \textbf{Habilidades, Certificaciones e Intereses}
\end{center}

\textbf{Certificaciones:} AWS Solutions Architect (2025), Azure AI Engineer Associate (AI-102) (2024, renovada 2025)

\textbf{Habilidades técnicas:} Ruby, Python, JavaScript, Java, SQL, Ruby on Rails, Node.js, React, TensorFlow, Docker, AWS, Azure, Git, LangChain, Redis, RabbitMQ

\textbf{Idiomas:} Español (Nativo), Inglés (Fluido)

\end{document}
